\documentclass[11pt, letterpaper]{article}

% --- Idioma y codificación ---
\usepackage[spanish]{babel}
\usepackage[utf8]{inputenc}
\usepackage[T1]{fontenc}
\usepackage{lmodern}

% --- Matemáticas ---
\usepackage{amsmath, amssymb, amsthm}

% --- Formato ---
\usepackage[margin=2.5cm]{geometry}
\usepackage{parskip}
\usepackage{booktabs}
\usepackage{listings}
\usepackage{hyperref}
\hypersetup{colorlinks=true, linkcolor=blue!70!black, urlcolor=blue!70!black}

% --- Código ---
\lstset{
  language=C, basicstyle=\ttfamily\small,
  keywordstyle=\bfseries, commentstyle=\itshape,
  numbers=left, numberstyle=\tiny, frame=single,
  breaklines=true
}

% --- Entornos de teoremas ---
\theoremstyle{definition}
\newtheorem{definition}{Definición}[section]
\theoremstyle{plain}
\newtheorem{theorem}{Teorema}[section]
\newtheorem{lemma}[theorem]{Lema}
\newtheorem{corollary}[theorem]{Corolario}
\theoremstyle{remark}
\newtheorem{remark}{Observación}[section]

\title{%
  T02 --- Demostración formal:\\[4pt]
  Correctitud y terminación de búsqueda binaria
}
\author{Curso Linux--C--Rust --- B15 Estructuras de datos en C}
\date{}

\begin{document}
\maketitle
\tableofcontents

% ===================================================================
\section{El algoritmo}
% ===================================================================

\begin{lstlisting}
int binary_search(int arr[], int n, int target) {
    int lo = 0;                              // S1
    int hi = n - 1;                          // S2
    while (lo <= hi) {                       // L1
        int mid = lo + (hi - lo) / 2;        // L2
        if (arr[mid] == target)               // L3
            return mid;                       // L4
        if (arr[mid] < target)                // L5
            lo = mid + 1;                     // L6
        else                                  // L7
            hi = mid - 1;                     // L8
    }
    return -1;                               // S3
}
\end{lstlisting}

\textbf{Precondición.} \texttt{arr[0..n)} está ordenado en orden no
decreciente:
$\texttt{arr}[0] \leq \texttt{arr}[1] \leq \ldots \leq \texttt{arr}[n{-}1]$.

\textbf{Postcondición.} El valor retornado es un índice $m$ tal que
$\texttt{arr}[m] = \texttt{target}$ (si existe), o $-1$ si
\texttt{target} no aparece en \texttt{arr[0..n)}.

% ===================================================================
\section{Demostración 1: correctitud (invariante de bucle)}
% ===================================================================

\begin{definition}[Invariante]
\[
  \textbf{INV}(\mathit{lo}, \mathit{hi}):
  \text{Si \texttt{target} existe en \texttt{arr[0..n)}, entonces}
  \;\exists\, k \text{ con } \mathit{lo} \leq k \leq \mathit{hi}
  \text{ tal que } \texttt{arr}[k] = \texttt{target}.
\]
\end{definition}

\begin{theorem}\label{thm:correctness}
\texttt{binary\_search} es correcta: satisface la postcondición.
\end{theorem}

\begin{proof}
Probamos el invariante en tres etapas.

\textbf{Inicialización} (después de S1--S2: $\mathit{lo} = 0$,
$\mathit{hi} = n - 1$). El rango $[\mathit{lo}, \mathit{hi}] = [0, n{-}1]$
es todo el array. Si \texttt{target} existe en el array, ciertamente
está en $\texttt{arr}[0..n{-}1]$. \checkmark

\textbf{Mantenimiento} (si el bucle continúa). Asumimos
$\text{INV}(\mathit{lo}, \mathit{hi})$ con $\mathit{lo} \leq \mathit{hi}$.
Se calcula $\mathit{mid} = \mathit{lo} +
\lfloor(\mathit{hi} - \mathit{lo})/2\rfloor$. Observación:
$\mathit{lo} \leq \mathit{mid} \leq \mathit{hi}$.

\emph{Caso A: $\texttt{arr}[\mathit{mid}] = \texttt{target}$} (L3--L4).
Se retorna $\mathit{mid}$. La postcondición se satisface directamente.
\checkmark

\emph{Caso B: $\texttt{arr}[\mathit{mid}] < \texttt{target}$} (L5--L6).
Se actualiza $\mathit{lo} = \mathit{mid} + 1$. Como el array está
ordenado:
\[
  \texttt{arr}[j] \leq \texttt{arr}[\mathit{mid}] < \texttt{target}
  \quad \forall\, j \leq \mathit{mid}
\]
Ningún elemento en $\texttt{arr}[0..\mathit{mid}]$ puede ser igual a
\texttt{target}. Si existe, debe estar en
$\texttt{arr}[\mathit{mid}{+}1..\mathit{hi}]$.
$\text{INV}(\mathit{mid}{+}1, \mathit{hi})$ vale. \checkmark

\emph{Caso C: $\texttt{arr}[\mathit{mid}] > \texttt{target}$} (L7--L8).
Se actualiza $\mathit{hi} = \mathit{mid} - 1$. Como el array está
ordenado:
\[
  \texttt{arr}[j] \geq \texttt{arr}[\mathit{mid}] > \texttt{target}
  \quad \forall\, j \geq \mathit{mid}
\]
Si \texttt{target} existe, debe estar en
$\texttt{arr}[\mathit{lo}..\mathit{mid}{-}1]$.
$\text{INV}(\mathit{lo}, \mathit{mid}{-}1)$ vale. \checkmark

\textbf{Terminación} ($\mathit{lo} > \mathit{hi}$). El rango
$[\mathit{lo}, \mathit{hi}]$ es vacío. Por el invariante, si
\texttt{target} existiera, estaría en ese rango vacío. Contradicción.
Retornar $-1$ es correcto. \checkmark
\end{proof}

% ===================================================================
\section{Demostración 2: terminación (función de decrecimiento)}
% ===================================================================

\begin{theorem}\label{thm:termination}
El bucle de búsqueda binaria termina en a lo sumo
$\lfloor \log_2 n \rfloor + 1$ iteraciones.
\end{theorem}

\begin{proof}
Definimos la función de decrecimiento (variant):
\[
  V = \mathit{hi} - \mathit{lo} + 1
\]
que es el tamaño del rango de búsqueda.

\textbf{Cota inferior.} El bucle solo ejecuta si
$\mathit{lo} \leq \mathit{hi}$, es decir, $V \geq 1$. Cuando el
bucle termina, $V \leq 0$.

\textbf{Decrecimiento estricto.} En cada iteración (que no retorna
en L4), $V$ decrece estrictamente.

\emph{Caso B} ($\mathit{lo} = \mathit{mid} + 1$):
\[
  V' = \mathit{hi} - (\mathit{mid} + 1) + 1
     = \mathit{hi} - \mathit{mid}
\]
Como $\mathit{mid} = \mathit{lo} +
\lfloor(\mathit{hi}{-}\mathit{lo})/2\rfloor$:
\[
  V' = \mathit{hi} - \mathit{lo}
     - \lfloor(\mathit{hi}{-}\mathit{lo})/2\rfloor
     = \lceil(\mathit{hi}{-}\mathit{lo})/2\rceil
     \leq \lfloor V/2 \rfloor
\]

\emph{Caso C} ($\mathit{hi} = \mathit{mid} - 1$):
\[
  V' = (\mathit{mid} - 1) - \mathit{lo} + 1
     = \mathit{mid} - \mathit{lo}
     = \lfloor(\mathit{hi}{-}\mathit{lo})/2\rfloor
     \leq \lfloor V/2 \rfloor
\]

\textbf{En ambos casos: $V' \leq \lfloor V/2 \rfloor$.} El rango se
reduce a la mitad (o menos) en cada iteración.

\textbf{Número de iteraciones.} Partiendo de $V_0 = n$ y reduciéndose
al menos a la mitad:
\[
  V_k \leq \left\lfloor \frac{n}{2^k} \right\rfloor
\]

El bucle termina cuando $V_k \leq 0$, lo que ocurre cuando
$\lfloor n/2^k \rfloor = 0$, es decir, $2^k > n$:
\[
  k > \log_2 n
  \quad\Longrightarrow\quad
  k \geq \lfloor \log_2 n \rfloor + 1 \qedhere
\]
\end{proof}

% ===================================================================
\section{Demostración 3: complejidad
  \texorpdfstring{$\Theta(\log n)$}{Theta(log n)} en el peor caso}
% ===================================================================

\begin{theorem}\label{thm:complexity}
La búsqueda binaria ejecuta $\Theta(\log n)$ comparaciones en el peor
caso.
\end{theorem}

\begin{proof}
\emph{Cota superior.} Demostrada en el Teorema~\ref{thm:termination}:
a lo sumo $\lfloor\log_2 n\rfloor + 1$ iteraciones, cada una con
$O(1)$ comparaciones.

\emph{Cota inferior.} Cualquier algoritmo de búsqueda basado en
comparaciones en un array ordenado de $n$ elementos necesita
$\Omega(\log n)$ comparaciones en el peor caso.

Argumento por árbol de decisión: cada comparación tiene 3 resultados
posibles ($<$, $=$, $>$), produciendo un árbol ternario. Hay $n$
resultados posibles (el elemento está en la posición $0, 1, \ldots,
n{-}1$) más 1 resultado ``no encontrado'', dando $n + 1$ hojas. Un
árbol ternario de altura~$h$ tiene a lo sumo $3^h$ hojas:
\[
  n + 1 \leq 3^h
  \quad\Longrightarrow\quad
  h \geq \log_3(n+1) = \Omega(\log n)
\]

Combinando: la búsqueda binaria es $\Theta(\log n)$ en el peor caso,
y por lo tanto es asintóticamente óptima.
\end{proof}

% ===================================================================
\section{Demostración 4: correctitud del cálculo de \texttt{mid}}
% ===================================================================

\begin{lemma}\label{lem:mid}
Si $0 \leq \mathit{lo} \leq \mathit{hi} < 2^{31}$, entonces
$\mathit{mid} = \mathit{lo} + (\mathit{hi} - \mathit{lo})/2$ no
produce overflow en aritmética de enteros de 32~bits con signo.
\end{lemma}

\begin{proof}
\begin{enumerate}
  \item $\mathit{hi} - \mathit{lo} \geq 0$ (porque
    $\mathit{lo} \leq \mathit{hi}$). No hay underflow.
  \item $\mathit{hi} - \mathit{lo} \leq \mathit{hi} < 2^{31}$. La
    diferencia cabe en un \texttt{int}.
  \item $(\mathit{hi} - \mathit{lo})/2 \leq \mathit{hi}/2 < 2^{30}$.
    La división reduce el valor.
  \item $\mathit{lo} + (\mathit{hi} - \mathit{lo})/2 \leq
    \mathit{lo} + (\mathit{hi} - \mathit{lo}) = \mathit{hi} < 2^{31}$.
    La suma no excede \texttt{hi}, que cabe en un \texttt{int}.
\end{enumerate}
\end{proof}

\begin{remark}[Contraste con la versión con bug]
$\mathit{mid} = (\mathit{lo} + \mathit{hi})/2$: si
$\mathit{lo} = \mathit{hi} = 2^{30}$, entonces
$\mathit{lo} + \mathit{hi} = 2^{31}$, que excede
$\texttt{INT\_MAX} = 2^{31} - 1$. Overflow con comportamiento
indefinido en C.
\end{remark}

% ===================================================================
\section{Resumen de resultados}
% ===================================================================

\begin{table}[h]
\centering
\begin{tabular}{llp{5.5cm}}
\toprule
\textbf{Resultado} & \textbf{Técnica} & \textbf{Clave} \\
\midrule
Correctitud
  & Invariante de bucle
  & Si target existe, está en \texttt{arr[lo..hi]} \\
Terminación
  & Función de decrecimiento
  & $V' \leq \lfloor V/2 \rfloor$ en cada iteración \\
Peor caso $\Theta(\log n)$
  & Cota superior + árbol de decisión
  & $\lfloor\log_2 n\rfloor + 1$ iter.; $\Omega(\log n)$ hojas \\
\texttt{mid} sin overflow
  & Análisis aritmético
  & $\mathit{lo} + (\mathit{hi}{-}\mathit{lo})/2 \leq \mathit{hi} < 2^{31}$ \\
\bottomrule
\end{tabular}
\end{table}

\end{document}
